
\subsection{博特--莫尔斯情形}
\label{subsec:3-7-5}

本小节把上一小节的结果推广到干净相交情形，也就是拉格朗日交弗洛尔理论的
博特--莫尔斯版本\liftIndex{I048}\liftIndex{I015}。（博特--莫尔斯理论在
弗洛尔同调中的早期工作见第 2.3.3 小节。）先考虑如下情形。令
$(L^{(0)},L^{(1)})$ 为一对相对 spin 拉格朗日子流形，并假设二者依下述意义
干净相交。

\begin{definition}{3.7.48}\label{def:3-7-48}
称 $L^{(0)}$ 与 $L^{(1)}$ \emph{干净相交}，是指如下条件成立。把交集分解成
连通分支
\[
 L^{(1)}\cap L^{(0)}
 =\coprod_{h\in\pi_0(L^{(0)}\cap L^{(1)})}R_h.
\]
假设每个 $R_h$ 都是光滑子流形，并且
\[
 N_{L^{(1)}/R_h}\oplus N_{L^{(0)}/R_h}\oplus TR_h
 =\bigl(TL^{(1)}+TL^{(0)}\bigr)|_{R_h},
\]
其中 $N_{L^{(0)}/R_h}$ 与 $N_{L^{(1)}/R_h}$ 分别是 $R_h$ 在
$L^{(0)}$ 与 $L^{(1)}$ 中的法丛。
\end{definition}

特别地，$L^{(0)}=L^{(1)}$ 是干净相交的特例。

记 $d_h=\dim R_h$。若 $L^{(1)}\cap L^{(0)}$ 不连通，就需要在每个分支上
固定一些附加数据。先在每个 $R_h$ 上固定基点 $p_h$，再考虑满足下列条件的
映射 $w\colon[0,1]\times[0,1]\to M$：
\begin{align}
 w(0,t)&=\ell_0(t),
 \tag{3.7.49.1}\label{eq:3-7-49-1}\\
 w(s,0)&\in L^{(0)},\qquad w(s,1)\in L^{(1)}
 \quad(0\leq s\leq1),
 \tag{3.7.49.2}\label{eq:3-7-49-2}\\
 w(1,t)&\equiv p_h.
 \tag{3.7.49.3}\label{eq:3-7-49-3}
\end{align}
这里 $\ell_0$ 是在道路空间
$\Omega(L^{(0)},L^{(1)})$ 的每个连通分支上固定的带基道路（见第 2.2 节）。
按定义 2.2.4 的方式在这些 $w$ 上定义等价关系，并以
$\widetilde\pi_0(L^{(0)}\cap L^{(1)})$ 表示全部二元组 $[h,w]$ 的等价类之集。

\begin{definition}{3.7.50}\label{def:3-7-50}
置
\[
 \widetilde\pi_0(L^{(0)}\cap L^{(1)})_{h_0}
 =
 \left\{
 [h,w]\in\widetilde\pi_0(L^{(0)}\cap L^{(1)})
 \ \middle|\ h=h_0
 \right\},
\]
并置
\[
 \Lambda_{h_0}
 =
 \operatorname{Span}_{\mathbb Q}
 \left\{
 [h,w_i]\ \middle|\
 [h,w_i]\in
 \widetilde\pi_0(L^{(0)}\cap L^{(1)})_{h_0}
 \right\}.
\]
\end{definition}

先回顾有限维博特--莫尔斯理论中为何必须计入临界子流形上的局部系统。设
$f\colon X\to\mathbb R$ 是有限维流形 $X$ 上的博特--莫尔斯函数，且其临界
点集分解为 $\operatorname{Cr}(f)=\coprod_hR_h$。这时存在谱序列
\[
 E_2^{ij}
 =
 \bigoplus_{\dim R_h=i}H^j(R_h;\Theta^-_{R_h})
 \Longrightarrow H^{i+j}(X;\mathbb Z).
\]
这里 $\Theta^-_{R_h}$ 是 $R_h$ 上不稳定丛的行列式实线丛。（这就是
\textbf{[Fuk96II]} 的定理 2.1；第一作者在那里漏写了 $\Theta^-_{R_h}$。
这一陈述原则上为拉乌尔·博特（Raoul Bott）所知，另见第 7.2.2 小节。）

在当前情形中，要以诺维科夫覆盖
$\widetilde\Omega(L^{(0)},L^{(1)})$ 代替 $X$，并以作用量泛函
$\mathcal A$ 代替 $f$，所以需要类似于上述 $\Theta^-_{R_h}$ 的数据。若把
$L^{(1)}$ 看成 $L^{(0)}$ 的余切丛的截面，所需局部系统正是有限维
博特--莫尔斯情形中的局部系统。一般地，定义 $R_h$ 上的向量丛
\begin{equation}
 V_h
 =
 \left.
 \frac{TL^{(1)}+TL^{(0)}}
 {(TL^{(1)}+TL^{(0)})^{\perp\omega}}
 \right|_{R_h}.
\tag{3.7.51}\label{eq:3-7-51}
\end{equation}
这是一个辛向量丛。以 $U,U'$ 分别表示 $TL^{(0)},TL^{(1)}$ 在 $V_h$ 中的
像；它们是 $V_h$ 的拉格朗日子丛。考虑 $V_h$ 中从 $U$ 到 $U'$ 的
拉格朗日子空间道路 $\lambda$，要求
$\lambda\oplus TR_h$ 是 $TM$ 中从 $T_pL^{(0)}$ 到 $T_pL^{(1)}$ 的
定向拉格朗日子空间道路。

接着考虑带拉格朗日边界条件的半圆盘上的多尔博算子之指标；第 8 章将作更详细
的讨论，见第 8.8 小节的命题 8.8.1。由此得到道路 $\lambda$ 的空间上的一族
算子。这一族算子之族指标的定向丛下降为 $R_h$ 上的局部系统，记作
$\Theta^-_{R_h}$。

定义 3.7.62 将对
$[h,w]\in\widetilde\pi_0(L^{(0)}\cap L^{(1)})$ 定义马斯洛夫指标
$\mu([h,w])\in\mathbb Z$。令
\[
 S^{*-\mu([h,w])}(R_h;\Theta^-_{R_h})
\]
为以局部系统 $\Theta^-_{R_h}\otimes\det TR_h$ 为系数的光滑奇异单形
上链复形。更明确地，若 $\Theta^-_{R_h}$ 由表示
$\rho_h\colon\pi_1(R_h)\to\{\pm1\}$ 诱导，则依附录 A2 的定义 A2.1 记
\[
 S^*(R_h;\Theta^-_{R_h})=S^*(R_h;\rho_h;\mathbb Q).
\]
注意 $R_h$ 在这里未必可定向；为了把链看作上链，需要因子 $\det TR_h$。
稍后将在命题 3.7.59 中从中取可数生成子复形
$C(R_h;\Theta^-_{R_h})$；我们假设它除次数平移外不依赖于 $w$。

定义 $\widehat{CF}_{BM}(L^{(1)},L^{(0)};\Lambda_{\mathrm{nov}})$ 为
\[
 \bigoplus_h
 \left(
 C(R_h;\Theta^-_{R_h})
 \otimes_{\mathbb Q}\Lambda_h
 \otimes_{\mathbb Q}\Lambda_{\mathrm{nov}}
 \right)
\]
的完备化，其中 $\Lambda_h$ 由定义 3.7.50 给出。若
\[
 S\otimes[h,w]\otimes T^\lambda e^\mu
 \in\widehat{CF}_{BM}(L^{(1)},L^{(0)};\Lambda_{\mathrm{nov}}),
 \qquad S\in C^\ell(R_h;\Theta^-_{R_h}),
\]
则把其能量定义为 $\int_w\omega+\lambda$，次数定义为
$k=\ell+\mu([h,w])+2\mu$。换言之，对 $R_h$ 中带标签 $[h,w]$ 的余维
$\ell$ 链 $S$，规定
\[
 \deg S=\ell+\mu([h,w]).
\]

仿照第 3.7.4 小节的 \eqref{eq:3-7-19-1}--\eqref{eq:3-7-19-3}，可以利用
能量和次数在
$\widehat{CF}_{BM}(L^{(1)},L^{(0)};\Lambda_{\mathrm{nov}})$ 上定义等价关系
$\sim$，也可以定义能量滤过；它在商空间上诱导滤过。

\begin{definition}{3.7.52}\label{def:3-7-52}
以 $C(L^{(1)},L^{(0)};\Lambda_{0,\mathrm{nov}})$ 表示
\[
 \widehat{CF}_{BM}(L^{(1)},L^{(0)};\Lambda_{\mathrm{nov}})/\sim
\]
关于能量滤过作完备化后的非负能量部分。
\end{definition}

赋予上述分次后，它成为滤过分次自由 $\Lambda_{0,\mathrm{nov}}$-模，并满足
(3.2.12.1)--(3.2.12.5)。

\begin{remark}{3.7.53}\label{rem:3-7-53}
若 $S=(|S|,f,\xi)$，其中 $f\colon|S|\to R_h$ 是单形的浸入，而 $\xi$ 是
$\det N_{R_h/S}\otimes\Theta^-_{R_h}$ 的平凡化，则得到
$C^*(L^{(1)},L^{(0)};\Lambda_{0,\mathrm{nov}})$ 中的元素 $[S,w]$。
这里对局部系统 $\Theta^-_{R_h}$ 的系数取庞加莱对偶。
\end{remark}

下面在 $C(L^{(1)},L^{(0)};\Lambda_{0,\mathrm{nov}})$ 上定义滤过
$A_\infty$ 双模结构；它的左、右作用代数分别为
\[
 \bigl(C(L^{(1)};\Lambda_{0,\mathrm{nov}}),\mathfrak m^{(1)}\bigr),
 \qquad
 \bigl(C(L^{(0)};\Lambda_{0,\mathrm{nov}}),\mathfrak m^{(0)}\bigr).
\]
为此先引入伪全纯映射的模空间。
令 $R_h,R_{h'}$ 是 $L^{(1)}\cap L^{(0)}$ 的连通分支，令
$\{J_t\}_t=\{J_t\}_{0\leq t\leq1}$ 是连接 $J_0,J_1$ 的单参数族，并取
$S\in C(R_h;\Theta^-_{R_h})$。

令 $\Sigma$ 为 $\mathbb R\times[0,1]$ 连同第 3.7.4 小节所述的球面气泡，
令 $u\colon\Sigma\to M$，并取 $\tau_j^{(i)}\in\mathbb R$（$i=0,1$）。
记
\[
 \vec\tau=(\vec\tau^{\,(0)},\vec\tau^{\,(1)}),
 \qquad
 \vec\tau^{\,(i)}
 =(\tau_1^{(i)},\ldots,\tau_{k_i}^{(i)}).
\]
考虑下列条件：
\begin{equation}
\begin{gathered}
 \frac{\partial u}{\partial\tau}
 +J_t\frac{\partial u}{\partial t}=0
 \quad\text{在 }\mathbb R\times[0,1]\text{ 上};\\
 \text{若球面气泡树扎根于 }(\tau,t),
 \text{ 则所有球面均为 }J_t\text{-全纯的}.
\end{gathered}
\tag{3.7.54.1}\label{eq:3-7-54-1}
\end{equation}
\begin{equation}
 u(\mathbb R\times\{0\})\subset L^{(0)},\qquad
 u(\mathbb R\times\{1\})\subset L^{(1)}.
\tag{3.7.54.2}\label{eq:3-7-54-2}
\end{equation}
\begin{equation}
\begin{gathered}
 u(\tau,t)\text{ 在 }\tau\to-\infty\text{ 时收敛到 }R_h\text{ 中一点},\\
 u(\tau,t)\text{ 在 }\tau\to+\infty\text{ 时收敛到 }R_{h'}\text{ 中一点}.
\end{gathered}
\tag{3.7.54.3}\label{eq:3-7-54-3}
\end{equation}
\begin{equation}
\begin{gathered}
 \tau_j^{(0)}<\tau_{j+1}^{(0)}
 \quad(j=1,\ldots,k_0-1),\\
 \tau_j^{(1)}>\tau_{j+1}^{(1)}
 \quad(j=1,\ldots,k_1-1).
\end{gathered}
\tag{3.7.54.4}\label{eq:3-7-54-4}
\end{equation}
\begin{equation}
 w\#u(\Sigma)\sim w',
\tag{3.7.54.5}\label{eq:3-7-54-5}
\end{equation}
其中 $\#$ 是显然的拼接，并且
\begin{equation}
 ((\Sigma,\vec\tau),u)\text{ 是稳定的}.
\tag{3.7.54.6}\label{eq:3-7-54-6}
\end{equation}

考虑由 $\tau$-平移给出的自然 $\mathbb R$-作用，以
$\sim_{\mathbb R}$ 表示相应等价关系。记
\[
\begin{aligned}
 &\mathring{\mathcal M}_{k_1,k_0}
 (L^{(1)},L^{(0)};[h,w],[h',w'];\{J_t\}_t)\\
 &\qquad
 =\mathring{\mathcal M}_{k_1,k_0}([h,w],[h',w'])\\
 &\qquad
 :=\left\{
 \begin{array}{c}
 \text{满足 \eqref{eq:3-7-54-1}--\eqref{eq:3-7-54-6} 的 }
 ((\Sigma,\vec\tau),u)\\
 \text{之全部 }\sim_{\mathbb R}\text{-等价类}
 \end{array}
 \right\}.
\end{aligned}
\]
它带有评价映射
\[
\begin{aligned}
 \operatorname{ev}
 &=(\operatorname{ev}^{(1)},\operatorname{ev}^{(0)})\\
 &=\bigl((\operatorname{ev}^{(1)}_1,\ldots,
 \operatorname{ev}^{(1)}_{k_1}),
 (\operatorname{ev}^{(0)}_1,\ldots,
 \operatorname{ev}^{(0)}_{k_0})\bigr)\\
 &\colon
 \mathring{\mathcal M}_{k_1,k_0}([h,w],[h',w'])
 \longrightarrow(L^{(1)})^{k_1}\times(L^{(0)})^{k_0},
\end{aligned}
\]
其中
\[
 \operatorname{ev}^{(i)}_j((\Sigma,\vec\tau),u)
 =u(\tau_j^{(i)},i)\in L^{(i)},\qquad i=0,1.
\]
还定义
\[
\begin{gathered}
 \operatorname{ev}_{-\infty}\colon
 \mathring{\mathcal M}_{k_1,k_0}([h,w],[h',w'])
 \longrightarrow R_h,\\
 \operatorname{ev}_{+\infty}\colon
 \mathring{\mathcal M}_{k_1,k_0}([h,w],[h',w'])
 \longrightarrow R_{h'},
\end{gathered}
\]
其值为
$\operatorname{ev}_{\pm\infty}((\Sigma,\vec\tau),u)
=\lim_{\tau\to\pm\infty}u(\tau,t)$。

\begin{proposition}{3.7.55}\label{prop:3-7-55}
模空间
\[
 \mathring{\mathcal M}_{k_1,k_0}
 ([h,w],[h',w'];\{J_t\}_t)
\]
可以紧化为
\[
 \mathcal M_{k_1,k_0}
 ([h,w],[h',w'];\{J_t\}_t).
\]
该紧化带有彼此相容、并同命题 3.4.2 中
\[
 \mathcal M_k^{\mathrm{main}}(L^{(0)};\beta^{(0)};J_0),
 \qquad
 \mathcal M_k^{\mathrm{main}}(L^{(1)};\beta^{(1)};J_1)
\]
上的仓西结构相容的仓西结构。此外，评价映射延拓为弱浸没
\[
\begin{aligned}
 (\operatorname{ev}_{-\infty},\operatorname{ev},
 \operatorname{ev}_{+\infty})\colon
 \mathcal M_{k_1,k_0}([h,w],[h',w'])
 \longrightarrow
 R_h\times(L^{(1)})^{k_1}\times(L^{(0)})^{k_0}\times R_{h'}.
\end{aligned}
\]
\end{proposition}

\begin{proof}
先仿照 \eqref{eq:3-7-29}，用
$\mathring{\mathcal M}_{k_1,k_0}([h,w],[h',w'])$ 定义
$\mathcal M^{\mathrm{unbr}}_{k_1,k_0}([h,w],[h',w'])$。
然后定义
\begin{equation}
\begin{aligned}
 &\mathcal M_{k_1,k_0}([h,w],[h',w'])\\
 &:=
 \coprod_{\substack{
 K,\ h(0),\ldots,h(K),\\
 w(0),\ldots,w(K),\
 k_{1,*},k_{0,*}}}
 \mathcal M^{\mathrm{unbr}}_{k_{1,1},k_{0,1}}
 ([h(0),w(0)],[h(1),w(1)])\\
 &\qquad{}_{R_{h(1)}}\!\times
 \mathcal M^{\mathrm{unbr}}_{k_{1,2},k_{0,2}}
 ([h(1),w(1)],[h(2),w(2)])
 {}_{R_{h(2)}}\!\times\cdots\\
 &\qquad\cdots{}_{R_{h(K-1)}}\!\times
 \mathcal M^{\mathrm{unbr}}_{k_{1,K},k_{0,K}}
 ([h(K-1),w(K-1)],[h(K),w(K)]).
\end{aligned}
\tag{3.7.56}\label{eq:3-7-56}
\end{equation}
右端不交并遍历所有满足下列条件的数据
\translatornote{原书 (3.7.56) 同时列出 $h(0),\ldots,h(K)$ 与
$k_{i,1},\ldots,k_{i,K}$，却只写从 $h(1)$ 到 $h(K)$ 的 $K-1$ 个因子；
(3.7.57.1) 又把 $h(1)$ 而非 $h(0)$ 绑定到输入。这会使 $h(0)$ 和
$k_{i,K}$ 均不进入构造。译文依上一小节的断裂链 \eqref{eq:3-7-31} 和本式
的纤维积端点，把它恢复为从 $(h(0),w(0))=(h,w)$ 到
$(h(K),w(K))=(h',w')$ 的 $K$ 段链。原书 (3.7.57.2) 还使用横截情形的
$\operatorname{Cr}(L^{(0)},L^{(1)})$；此处端点类已由定义 3.7.50 改为
$\widetilde\pi_0(L^{(0)}\cap L^{(1)})$，译文一并恢复型别。}：
\begin{align}
 h(0)&=h,\quad w(0)=w,\qquad
 h(K)=h',\quad w(K)=w',
\tag{3.7.57.1}\label{eq:3-7-57-1}\\
 [h(i),w(i)]&\in
 \widetilde\pi_0(L^{(0)}\cap L^{(1)})
 \quad(0\leq i\leq K),
\tag{3.7.57.2}\label{eq:3-7-57-2}\\
 k_{1,1}+\cdots+k_{1,K}&=k_1,\qquad
 k_{0,1}+\cdots+k_{0,K}=k_0.
\tag{3.7.57.3}\label{eq:3-7-57-3}
\end{align}
注意 \eqref{eq:3-7-56} 在各 $R_{h(i)}$ 上取纤维积；横截情形则取直接积，
见 \eqref{eq:3-7-31}
\translatornote{原书此处误引 (3.7.33)，该式只是评价映射；直接积定义实际
是 (3.7.31)。}。仿照第 7.1 节，可以在
$\mathcal M_{k_1,k_0}([h,w],[h',w'])$ 上定义拓扑并证明它是豪斯多夫紧空间；
仓西结构的构造也与第 7.1 节相同。
\end{proof}

即使 $(L^{(0)},L^{(1)})$ 是相对 spin 对，模空间
$\mathcal M_{k_1,k_0}([h,w],[h',w'])$ 也未必可定向。关于这一点的精确
陈述见下文引理 3.7.70；博特--莫尔斯情形的定向问题将在第 8.8 节详述。

下面考虑 $S(L^{(0)};\mathbb Q),S(L^{(1)};\mathbb Q)$ 的子模
$C(L^{(0)};\mathbb Q),C(L^{(1)};\mathbb Q)$；它们分别由第 3.5 节的
可数集 $\mathcal X_1(L^{(0)}),\mathcal X_1(L^{(1)})$ 生成。
令
\[
 O_{R_h}=\det TR_h
\]
为 $R_h$ 的切丛的定向实线丛。在下一个命题的证明中，还要选取光滑链
$S=(|S|,f)$ 的可数集 $\mathcal X_1(R_h)$，并给每条链配备
$f^*O_{R_h}\otimes f^*\Theta^-_{R_h}$ 的平凡化。
$S\in\mathcal X_1(R_h)$ 表示 $R_h$ 中取值于 $\Theta^-_{R_h}$ 的上链，
仍以 $S$ 表示。它们在 $S(R_h;\Theta^-_{R_h})$ 中生成子模
$C(R_h;\Theta^-_{R_h})$。

\begin{remark}{3.7.58}\label{rem:3-7-58}
定义滤过 $A_\infty$ 代数时，光滑奇异单形的可数集
$\mathcal X_1(L^{(i)})$ 已经选定。在命题 3.7.59 的证明中，我们另选
$\mathcal X_1(R_h)$ 来定义 $C(R_h;\Theta^-_{R_h})$；论证与第 7.2 节类似。
\end{remark}

令 $(P_j^{(i)},f_j^{(i)})\in\mathcal X_1(L^{(i)})$，
$j=1,\ldots,k_i$，并令
$S=(|S|,f)\in\mathcal X_1(R_h)$。取纤维积
\[
\begin{aligned}
 &\mathcal M_{k_1,k_0}
 (L^{(1)},L^{(0)};[h,w],[h',w'])\\
 &\quad{}_{(\operatorname{ev}^{(1)},\operatorname{ev}_{-\infty},
 \operatorname{ev}^{(0)})}
 \mathop{\times}_{\prod_j f_j^{(1)}\times f\times\prod_j f_j^{(0)}}
 \left(
 \prod_{j=1}^{k_1}P_j^{(1)}
 \times S\times
 \prod_{j=1}^{k_0}P_j^{(0)}
 \right),
\end{aligned}
\]
并把它记作
\[
\begin{aligned}
 \mathcal M_{k_1,k_0,(+\infty)}
 (L^{(1)},L^{(0)};[h,w],[h',w'];
 \vec P^{\,(1)},S,\vec P^{\,(0)}),
\end{aligned}
\]
其中
$\vec P^{\,(i)}=(P_1^{(i)},\ldots,P_{k_i}^{(i)})$。
由命题 3.7.55，这一空间带有仓西结构；而且评价映射
\[
 \operatorname{ev}_{+\infty}\colon
 \mathcal M_{k_1,k_0,(+\infty)}
 (L^{(1)},L^{(0)};[h,w],[h',w'];
 \vec P^{\,(1)},S,\vec P^{\,(0)})
 \longrightarrow R_{h'}
\]
是弱浸没。它的维数照例由适当边界条件下的柯西--黎曼型算子之指标给出；
该算子来自 \eqref{eq:3-7-54-1} 的线性化。稍后再给出这个指标的拓扑定义，
先陈述下述命题。

\liftIndexRangeStart{I129}
\begin{proposition}{3.7.59}\label{prop:3-7-59}
在上述情形中，进一步假设 $(L^{(0)},L^{(1)})$ 是相对 spin 对。则对每个
$h\in\pi_0(L^{(0)}\cap L^{(1)})$，都存在光滑奇异单形的可数集
$\mathcal X_1(R_h)$，并且可为每个 $[h,w]$ 指定整数
$\mu([h,w])\in\mathbb Z$，使得在
\[
 \mathcal M_{k_1,k_0,(+\infty)}
 (L^{(1)},L^{(0)};[h,w],[h',w'];
 \vec P^{\,(1)},S,\vec P^{\,(0)})
\]
上可以找到一族多值截面
$\mathfrak s_{[h,w],[h',w'],\vec P^{\,(1)},S,\vec P^{\,(0)}}$，
满足下列性质：
\begin{equation}
\begin{aligned}
 &\dim
 \mathcal M_{k_1,k_0,(+\infty)}
 (L^{(1)},L^{(0)};[h,w],[h',w'];
 \vec P^{\,(1)},S,\vec P^{\,(0)})\\
 &\quad=
 \dim S+d_{h'}-d_h+\mu([h',w'])-\mu([h,w])\\
 &\qquad\quad
 -\sum_{j=1}^{k_0}\deg P_j^{(0)}
 -\sum_{j=1}^{k_1}\deg P_j^{(1)}
 +k_0+k_1-1.
\end{aligned}
\tag{3.7.60.1}\label{eq:3-7-60-1}
\end{equation}
这里 $d_h=\dim R_h$，而 $\deg P_j^{(0)},\deg P_j^{(1)}$ 是移位前作为
上链的次数。
\begin{equation}
 \mathfrak s_{[h,w],[h',w'],\vec P^{\,(1)},S,\vec P^{\,(0)}}
 \text{ 横截于零}.
\tag{3.7.60.2}\label{eq:3-7-60-2}
\end{equation}
\begin{equation}
\begin{gathered}
 \mathfrak s_{[h,w],[h',w'],\vec P^{\,(1)},S,\vec P^{\,(0)}}
 \text{ 与其余同类多值截面相容，}\\
 \text{并与 }
 \mathfrak s_{\beta^{(0)},\vec P^{\,(0)}},
 \mathfrak s_{\beta^{(1)},\vec P^{\,(1)}}
 \text{ 相容}.
\end{gathered}
\tag{3.7.60.3}\label{eq:3-7-60-3}
\end{equation}
\begin{equation}
 \left(
 \mathfrak s_{[h,w],[h',w'],\vec P^{\,(1)},S,\vec P^{\,(0)}}^{-1}(0),
 \operatorname{ev}_{+\infty}
 \right)
 \in C(R_{h'};\Theta^-_{R_{h'}}).
\tag{3.7.60.4}\label{eq:3-7-60-4}
\end{equation}
\end{proposition}

要使 \eqref{eq:3-7-60-4} 完全严格，还须如 \eqref{eq:3-5-3-4} 之后那样，
选取三角剖分及其顶点次序，从而把虚基本链视作奇异链。这里的处理完全相同，
不再重复。

这个命题的证明与命题 3.7.36 相似，不过有三个新问题。第一个是
$\mathcal X_1(R_h)$ 的选择，可按第 7.2 节的方法处理；另两个分别是次数
$\mu([h,w])$ 与定向平坦丛 $\Theta^-_{R_h}$。下面说明后两点。

为给出维数公式，先描述 \eqref{eq:3-7-54-1} 的线性化问题。回忆第 2.2.2
小节已经在道路空间 $\Omega(L^{(0)},L^{(1)})$ 的每个连通分支上固定基道路
$\ell_0\colon[0,1]\to M$，其中
$\ell_0(0)\in L^{(0)}$、$\ell_0(1)\in L^{(1)}$
\translatornote{原书括注把两个端点都写在 $L^{(0)}$ 中；紧随其后的
$\lambda_0(1)=T_{\ell_0(1)}L^{(1)}$ 以及道路空间
$\Omega(L^{(0)},L^{(1)})$ 均表明第二处应为 $L^{(1)}$。}。
还选取了道路 $t\mapsto\lambda_0(t)$，其中 $\lambda_0(t)$ 是
$T_{\ell_0(t)}M$ 的定向拉格朗日线性子空间，并且作为定向线性空间有
\[
 \lambda_0(0)=T_{\ell_0(0)}L^{(0)},\qquad
 \lambda_0(1)=T_{\ell_0(1)}L^{(1)}.
\]

对 $p_h\in R_h\subset L^{(0)}\cap L^{(1)}$，考虑拉格朗日子空间的道路空间
\[
\begin{aligned}
 \mathcal P_{R_h}(T_{p_h}L^{(0)},T_{p_h}L^{(1)})
 :=\{\,\lambda\colon[0,1]\to\Lambda_{\mathrm{ori}}(T_{p_h}M)
 \mid{}&
 \lambda(0)=T_{p_h}L^{(0)},\\[-2pt]
 &\lambda(1)=T_{p_h}L^{(1)}\,\}.
\end{aligned}
\]
这里 $\Lambda_{\mathrm{ori}}(T_{p_h}M)$ 是 $T_{p_h}M$ 的定向拉格朗日
格拉斯曼流形。

利用 $\lambda_0$ 与 \eqref{eq:3-7-49-1}--\eqref{eq:3-7-49-3} 中的 $w$，
如下确定 $\mathcal P_{R_h}(T_{p_h}L^{(0)},T_{p_h}L^{(1)})$ 中一个元素的
同伦类。因为 $[0,1]^2$ 可缩，所以有平凡化
\[
 w^*TM\cong[0,1]^2\times T_{p_h}M.
\]
在
\[
 (\{0\}\times[0,1])\cup([0,1]\times\{0,1\})
 \cong[0,1]
\]
上定义
\[
 \lambda(0,t)=\lambda_0(t),\qquad
 \lambda(s,i)=T_{w(s,i)}L^{(i)}\quad(i=0,1).
\]
所得道路记为 $\lambda_{w,\lambda_0}$。

\begin{figure}[htbp]
\centering
\setlength{\unitlength}{1mm}
\begin{picture}(94,57)
  \put(25,8){\line(0,1){40}}
  \put(25,48){\line(1,0){45}}
  \put(25,8){\line(1,0){45}}
  \multiput(70,8)(0,4){10}{\line(0,1){2}}
  \put(25,28){\circle*{1.5}}
  \put(47.5,48){\circle*{1.5}}
  \put(47.5,8){\circle*{1.5}}
  \put(7,26){$\lambda_0(t)$}
  \put(39,51){$T_{w(s,1)}L^{(1)}$}
  \put(39,1){$T_{w(s,0)}L^{(0)}$}
\end{picture}
  \par\smallskip
  \renewcommand{\thefigure}{3.7.8}
  \refstepcounter{figure}
  \textbf{\figurename~\thefigure}\label{fig:3-7-8}
\end{figure}
\liftIndexRangeEnd{I129}

置
\[
\begin{aligned}
 Z_-&=\{z\in\mathbb C\mid |z|\leq1\}
 \cup\{z\in\mathbb C\mid\operatorname{Re}z\geq0,\ 
 |\operatorname{Im}z|\leq1\},\\
 Z_+&=\{z\in\mathbb C\mid\operatorname{Re}z\leq0,\ 
 |\operatorname{Im}z|\leq1\}
 \cup\{z\in\mathbb C\mid |z|\leq1\}.
\end{aligned}
\]

\begin{figure}[htbp]
\centering
\setlength{\unitlength}{1mm}
\begin{picture}(104,73)
  \put(42,57){\line(1,0){48}}
  \put(42,45){\line(1,0){48}}
  \qbezier(42,45)(30,45)(30,51)
  \qbezier(30,51)(30,57)(42,57)
  \put(42,39){\line(0,1){24}}
  \put(39.7,45){\line(1,1){12}}
  \put(52,45){\line(1,1){12}}
  \put(63.3,45){\line(1,1){12}}
  \put(74.4,45){\line(1,1){12}}
  \put(85.5,45){\line(1,1){4.5}}
  \put(30,51){\circle*{1.4}}
  \put(20,49){$-1$}
  \put(44,59){$\sqrt{-1}$}
  \put(44,39){$-\sqrt{-1}$}
  \put(31,34){$Z_-$}
  \put(9,24){\line(1,0){48}}
  \put(9,12){\line(1,0){48}}
  \qbezier(57,24)(69,24)(69,18)
  \qbezier(69,18)(69,12)(57,12)
  \put(57,6){\line(0,1){24}}
  \put(10,24){\line(1,-1){12}}
  \put(21.2,24){\line(1,-1){12}}
  \put(32.7,24){\line(1,-1){12}}
  \put(44.1,24){\line(1,-1){12}}
  \put(55.6,24){\line(1,-1){7.4}}
  \put(69,18){\circle*{1.4}}
  \put(72,16){$1$}
  \put(59,26){$\sqrt{-1}$}
  \put(59,6){$-\sqrt{-1}$}
  \put(45,0){$Z_+$}
\end{picture}
  \par\smallskip
  \renewcommand{\thefigure}{3.7.9}
  \refstepcounter{figure}
  \textbf{\figurename~\thefigure}\label{fig:3-7-9}
\end{figure}

记 $z=\tau+\sqrt{-1}\,t$。对每个
$\lambda\in\mathcal P_{R_h}(T_{p_h}L^{(0)},T_{p_h}L^{(1)})$，如下定义一个
弗雷德霍姆算子。令
$W_\lambda^{1,p}(Z_+;T_{p_h}M)$（分别地，
$W_\lambda^{1,p}(Z_-;T_{p_h}M)$）为由局部 $L^{1,p}$ 映射
$\zeta_+\colon Z_+\to T_{p_h}M$（分别为
$\zeta_-\colon Z_-\to T_{p_h}M$）组成的巴拿赫空间，并要求：
\begin{align}
 \zeta_\pm(\tau,i)&\in T_{p_h}L^{(i)}
 \quad(i=0,1),
\tag{3.7.61.1}\label{eq:3-7-61-1}\\
 \zeta_\pm(z_\pm(t))&\in\lambda(t),\qquad
 z_\pm(t)=e^{\pi\sqrt{-1}(-1/2\pm t)}\in\partial Z_\pm,
\tag{3.7.61.2}\label{eq:3-7-61-2}\\
 \int_{Z_\pm}e^{\delta|\tau|}
 \bigl(|\nabla\zeta_\pm|^p+|\zeta_\pm|^p\bigr)\,d\tau\,dt
 &<\infty.
\tag{3.7.61.3}\label{eq:3-7-61-3}
\end{align}
令
\[
 L^p(Z_\pm;T_{p_h}M\otimes\Lambda^{0,1}(Z_\pm))
\]
为局部 $L^p$ 截面 $\eta_\pm$ 组成的巴拿赫空间，并要求
\[
 \int_{Z_\pm}e^{\delta|\tau|}|\eta_\pm|^p\,d\tau\,dt<\infty.
\]
多尔博算子诱导有界线性映射
\[
 \bar\partial_{\lambda,Z_\pm}\colon
 W_\lambda^{1,p}(Z_\pm;T_{p_h}M)
 \longrightarrow
 L^p(Z_\pm;T_{p_h}M\otimes\Lambda^{0,1}(Z_\pm)).
\]
众所周知，$\bar\partial_{\lambda,Z_\pm}$ 是弗雷德霍姆算子。

\begin{definition}{3.7.62}\label{def:3-7-62}
记
\[
 \mu_{\lambda_0}^{\mathrm{anal}}([h,w])
 =
 \operatorname{Index}
 \bar\partial_{\lambda_{w,\lambda_0},Z_-}.
\]
\end{definition}

引理 3.7.69 将证明，这个解析指标等于广义马斯洛夫--莫尔斯指标
$\mu([h,w];\lambda_0)$。

\par\noindent\textit{命题 3.7.59 的证明。}\quad
置
$\mu([h,w])=\mu_{\lambda_0}^{\mathrm{anal}}([h,w])$。
先验证 \eqref{eq:3-7-60-1}。取
\[
 u\in\mathcal M(L^{(1)},L^{(0)};[h,w],[h',w']),
\]
即 $u\colon\mathbb R\times[0,1]\to M$ 是伪全纯映射。考虑巴拿赫空间
\[
 W_\delta^{1,p}
 (\mathbb R\times[0,1];u^*TM,u^*TL^{(1)},u^*TL^{(0)}),
\]
其元素为 $u^*TM$ 的局部 $L^{1,p}$ 截面 $\zeta$，并满足
\translatornote{原书 (3.7.63.1) 把变化点处的边界条件写成
$\zeta(\tau,i)\in T_{p_h}L^{(i)}$；但 $\zeta$ 是 $u^*TM$ 的截面，且
$u(\tau,i)$ 随 $\tau$ 变化，故其正确纤维只能是
$T_{u(\tau,i)}L^{(i)}$。}：
\begin{align}
 \zeta(\tau,i)&\in T_{u(\tau,i)}L^{(i)}
 \quad(i=0,1),
\tag{3.7.63.1}\label{eq:3-7-63-1}\\
 \int_{\mathbb R\times[0,1]}e^{\delta|\tau|}
 \bigl(|\nabla\zeta|^p+|\zeta|^p\bigr)\,d\tau\,dt
 &<\infty.
\tag{3.7.63.2}\label{eq:3-7-63-2}
\end{align}
再令
\[
 L^p(\mathbb R\times[0,1];
 u^*TM\otimes\Lambda^{0,1}(\mathbb R\times[0,1]))
\]
为满足
\[
 \int_{\mathbb R\times[0,1]}e^{\delta|\tau|}|\eta|^p\,d\tau\,dt<\infty
\]
的局部 $L^p$ 截面 $\eta$ 所成的巴拿赫空间。多尔博算子诱导有界线性映射
\[
\begin{aligned}
 \bar\partial_u\colon&
 W_\delta^{1,p}
 (\mathbb R\times[0,1];u^*TM,u^*TL^{(1)},u^*TL^{(0)})\\
 &\longrightarrow
 L^p(\mathbb R\times[0,1];
 u^*TM\otimes\Lambda^{0,1}(\mathbb R\times[0,1])).
\end{aligned}
\]
因为 $w\#u\sim w'$，标准的指标求和公式给出下述引理
\translatornote{原书把下述式子的 $+d_h$ 排在等号右边；但引理 3.7.65
给出
$\dim\mathcal M=\operatorname{Index}\bar\partial_u+d_h+d_{h'}-1$，
而本页紧接着明确由两引理推出
$\dim\mathcal M=\mu([h',w'])-\mu([h,w])+d_{h'}-1$。三者相容要求
$+d_h$ 位于下式左边；引理证明中 $-\infty$ 端的 $d_h$ 维零特征空间也支持
这一指标拼接式。译文据此恢复。}。

\begin{lemma}{3.7.64}\label{lem:3-7-64}
\[
 \operatorname{Index}\bar\partial_{\lambda_{w,\lambda_0},Z_-}
 +\operatorname{Index}\bar\partial_u+d_h
 =
 \operatorname{Index}\bar\partial_{\lambda_{w',\lambda_0},Z_-}.
\]
\end{lemma}

\begin{proof}
当 $\tau\to+\infty$ 时，
$\bar\partial_{\lambda_{w,\lambda_0},Z_-}$ 是乘积型的，即形如
$\partial/\partial\tau+P_\tau$；而
$\lim_{\tau\to+\infty}P_\tau$ 的零特征值个数恰为
$d_h=\dim R_h$。算子 $\bar\partial_u$ 在 $\tau\to-\infty$ 时具有类似
形式，零特征值个数也是 $d_h$。对 \textbf{[APS75]} 的定理 3.10 作显然
修改，便由这些事实得到结论。
\end{proof}

\begin{lemma}{3.7.65}\label{lem:3-7-65}
\[
 \dim\mathcal M(L^{(1)},L^{(0)};[h,w],[h',w'])
 =
 \operatorname{Index}\bar\partial_u+d_h+d_{h'}-1.
\]
\end{lemma}

\begin{proof}
$\bar\partial_u$ 与非线性柯西--黎曼方程在 $u$ 处的线性化算子
$D_u\bar\partial$ 具有相同主符号。由 \eqref{eq:3-7-63-2}，截面 $\zeta$
在 $|\tau|\to\infty$ 时指数衰减。因此
$\operatorname{Index}\bar\partial_u$ 是收敛到固定点
$p_h\in R_h$、$p_{h'}\in R_{h'}$ 的伪全纯条带模空间之虚维数。
另一方面，$\mathcal M(L^{(1)},L^{(0)};[h,w],[h',w'])$ 中的映射只要求在
$\tau\to-\infty,+\infty$ 时分别收敛到 $R_h,R_{h'}$ 中的某些点，并且还要
除以 $\mathbb R$-作用，故维数正如所述。$R_h=L^{(0)}=L^{(1)}$ 时的
弗雷德霍姆理论见第 7.1.2 小节；推广到这里任意干净相交对是直接的。
\end{proof}

继续证明 \eqref{eq:3-7-60-1}。空间
\[
 \mathcal M_{k_1,k_0,(+\infty)}
 (L^{(1)},L^{(0)};[h,w],[h',w'];
 \vec P^{\,(1)},S,\vec P^{\,(0)})
\]
由 $\mathcal M(L^{(1)},L^{(0)};[h,w],[h',w'])$ 与若干链取纤维积得到。
同 $S$ 取纤维积使维数减少 $d_h-\dim S$；同 $P_j^{(i)}$ 取纤维积使维数
减少 $\deg P_j^{(i)}$。加入 $k_1+k_0$ 个边界点则增加 $k_1+k_0$ 维。
所以
\[
\begin{aligned}
 &\dim\mathcal M_{k_1,k_0,(+\infty)}
 (L^{(1)},L^{(0)};[h,w],[h',w'];
 \vec P^{\,(1)},S,\vec P^{\,(0)})\\
 &\quad=
 \dim\mathcal M(L^{(1)},L^{(0)};[h,w],[h',w'])
 +\dim S-d_h\\
 &\qquad
 -\sum_{j=1}^{k_0}\deg P_j^{(0)}
 -\sum_{j=1}^{k_1}\deg P_j^{(1)}
 +k_0+k_1.
\end{aligned}
\]
引理 3.7.64 与引理 3.7.65 蕴含
\[
 \dim\mathcal M(L^{(1)},L^{(0)};[h,w],[h',w'])
 =
 \mu_{\lambda_0}([h',w'])-\mu_{\lambda_0}([h,w])+d_{h'}-1.
\]
于是 \eqref{eq:3-7-60-1} 成立。这里所有维数都是虚维数，即依定义为仓西
结构意义下的维数。

\begin{remark}{3.7.66}\label{rem:3-7-66}
公式 \eqref{eq:3-7-60-1} 只把整数 $\mu([h,w])$ 确定到整体加一个常数。
这些整数依赖于道路空间 $\Omega(L^{(0)},L^{(1)})$ 的连通分支；在上述讨论
中，该分支由 $\ell_0$ 确定。只要把它们定义成满足维数公式
\eqref{eq:3-7-60-1}，这一常数歧义就不会造成问题。

不过，为构造双模，特别是研究有关模空间的定向，必须使用
$\mu([h,w])$ 的一个典范选择，也就是仅依赖于 $\lambda_0$ 的绝对次数。
定义 3.7.62 的解析指标
$\mu_{\lambda_0}^{\mathrm{anal}}([h,w])$ 就提供了这样的选择。

例如，在公式 (3.7.2) 中，$\deg y$（在当前几何情形中即
$\deg S$，见 (8.8.8)）出现在右端第三项，却不出现在第二项。若把
$\mu([h,w])$ 整体平移一个常数，则 $\deg S=\deg y$ 随之改变，而第二项中
的 $\deg x_{i,j}=\deg P_j^{(i)}$ 不变。因此除非适当归一化
$\mu([h,w])$，否则 (3.7.2) 不再成立。事实上，第 8 章赋予
$\mathcal M(L^{(1)},L^{(0)};[h,w],[h',w'])$ 的定向将依赖参考拉格朗日道路
$\lambda_0$，所以运算 $\mathfrak n$ 中的符号也依赖于 $\lambda_0$。
\end{remark}

在横截拉格朗日子流形对的情形，第 2.2 节已经指出解析指标等于拓扑指标
$\mu([h,w];\lambda_0)$（定义 2.2.12），后者称为马斯洛夫--莫尔斯指标。
下面把它推广到博特--莫尔斯情形，并证明两种指标仍然相等。

考虑命题 3.7.59 的情形，其中 $L^{(0)}\cap L^{(1)}$ 是干净相交。
固定 $h\in\pi_0(L^{(0)}\cap L^{(1)})$，并取代表点
$p_h\in L^{(0)}\cap L^{(1)}$。令
$w\colon[0,1]^2\to M$、$\ell_0$ 与 $\lambda_0$ 为满足
\eqref{eq:3-7-49-1}--\eqref{eq:3-7-49-3} 的光滑数据。对
$(s,0),(s,1)\in\partial[0,1]^2$，记
\[
 \lambda_w(s,i)=T_{w(s,i)}L^{(i)},\qquad i=0,1.
\]
现在把这些拉格朗日道路补成一个丛对。为此，需要选取
$\Lambda(T_{p_h}M)$ 中的拉格朗日道路 $\alpha_h$，满足
\begin{equation}
 \alpha_h(0)=T_{p_h}L^{(0)},\qquad
 \alpha_h(1)=T_{p_h}L^{(1)}.
\tag{3.7.67.1}\label{eq:3-7-67-1}
\end{equation}
选择 $\alpha_h(t)$，使它是 $T_{p_h}M$ 中定向拉格朗日子空间的道路，并
延拓 \eqref{eq:3-7-67-1} 两端给定的定向。与第 2.1 节一样，这类道路的
同伦类并不唯一，因此还要用类似 \eqref{eq:2-1-4-3} 的条件固定其同伦类。
具体地，要求
\begin{equation}
 \alpha_h(t)\cap T_{p_h}L^{(1)}
 =
 T_{p_h}L^{(0)}\cap T_{p_h}L^{(1)}
 \qquad(0\leq t<1).
\tag{3.7.67.2}\label{eq:3-7-67-2}
\end{equation}
把
\[
 T_{p_h}M=\mathbb C^n
 =\mathbb C^{d_h}\oplus\mathbb C^{n-d_h},
 \qquad
 \mathbb C^{d_h}=T_{p_h}R_h\otimes\mathbb C,
\]
作此识别后，可以把 $\alpha_h(t)$ 分解为
\begin{equation}
 \alpha_h(t)=T_{p_h}R_h\oplus\alpha_h^2(t).
\tag{3.7.67.3}\label{eq:3-7-67-3}
\end{equation}
并要求在充分小的 $t$ 上
\begin{equation}
 \alpha_h^2(t)\in U_0(\alpha_h^2(0)),
\tag{3.7.67.4}\label{eq:3-7-67-4}
\end{equation}
其中右端是在 $\mathbb C^{n-d_h}$ 上应用 (2.1.2) 所得的集合。

于是得到辛丛对
$(w^*TM,\widetilde\lambda_{0,\lambda_0,h})$（定义 2.1.13）；这里
$\widetilde\lambda_{0,\lambda_0,h}$ 在 $(0,t)$、$(1,t)$ 与 $(s,i)$ 上
分别取 $\lambda_0(t)$、$\alpha_h(t)$ 与 $\lambda_w(s,i)$。定义
\begin{equation}
 \mu([h,w];\lambda_0)
 :=
 \mu(w^*TM,\widetilde\lambda_{0,\lambda_0,h}).
\tag{3.7.68}\label{eq:3-7-68}
\end{equation}
\liftIndex{I023}\liftIndex{I056}
当 $L^{(0)}$ 与 $L^{(1)}$ 横截时，这一定义退化为定义 2.2.12。

\begin{lemma}{3.7.69}\label{lem:3-7-69}
设 $w$ 如上，而
$\bar\partial_{\lambda_{w,\lambda_0},Z_-}$ 是满足
\eqref{eq:3-7-61-1}--\eqref{eq:3-7-61-3} 的多尔博算子，则
\[
 \mu([h,w];\lambda_0)
 =
 \mu_{\lambda_0}^{\mathrm{anal}}([h,w])
 =
 \operatorname{Index}\bar\partial_{\lambda_{w,\lambda_0},Z_-}.
\]
\end{lemma}

\begin{proof}
在平凡化
$w^*TM\cong[0,1]^2\times T_{p_h}M$ 下，拉格朗日子丛
$\widetilde\lambda_{0,\lambda_0,h}$ 给出一条拉格朗日回路，仍以同一符号
表示。回路的马斯洛夫指标在辛变换和同伦下不变。考虑固定边界段
$\{(1,t)\mid t\in[0,1]\}\subset\partial[0,1]^2$ 的同伦，再用第 2.1 节的
截断函数 $\chi\colon[0,\infty)\to[0,1]$ 把
$[0,1)\times[0,1]\cong\mathbb R\times[0,1]$ 拉长。在这一拉长下，
$\bar\partial_{\lambda_{w,\lambda_0},Z_-}$ 在 $\tau\to\infty$ 时成为
乘积型；由弗雷德霍姆指标的同伦不变性，只要保持渐近条件，其指标在系数和
边界条件的连续形变下不变。

因此只须考虑下述特殊情形；这类约化的详细说明见 \textbf{[Oh99]} 的附录：
\[
\begin{gathered}
 T_{p_h}M=\mathbb C^n,\qquad
 T_{p_h}L^{(0)}=\mathbb R^n,\qquad
 T_{p_h}L^{(1)}
 =\mathbb R^{d_h}\oplus\sqrt{-1}\,\mathbb R^{n-d_h},\\
 \alpha_h(t)
 =\mathbb R^{d_h}\oplus
 e^{\pi\sqrt{-1}t/2}\mathbb R^{n-d_h},
\end{gathered}
\]
而
$\lambda_{w,\lambda_0}\colon
\partial[0,1]^2\setminus\{(1,t)\}\to\mathbb C^n$ 为
\[
\begin{aligned}
 \lambda_{w,\lambda_0}(s,0)&\equiv\mathbb R^n,\\
 \lambda_{w,\lambda_0}(s,1)&\equiv
 \mathbb R^{d_h}\oplus\sqrt{-1}\,\mathbb R^{n-d_h},\\
 \lambda_{w,\lambda_0}(0,t)&=
 \mathbb R^{d_h}\oplus
 \left(
 e^{-(l_1+1/2)\pi\sqrt{-1}t}\mathbb R
 \oplus\cdots\oplus
 e^{-(l_{n-d_h}+1/2)\pi\sqrt{-1}t}\mathbb R
 \right),
\end{aligned}
\]
其中 $l_j\in\mathbb Z$。于是拉格朗日回路为
$\widetilde\lambda_{0,\lambda_0,h}
=\lambda_{w,\lambda_0}\cup\alpha_h$。

现在指标问题分裂成若干一维问题：
\[
 \begin{cases}
  \bar\partial\zeta=0,\\
  \zeta(\tau,0)\in\mathbb R,\quad
  \zeta(\tau,1)\in\sqrt{-1}\,\mathbb R,\\
  \zeta(0,t)\in
  e^{-(l+1/2)\pi\sqrt{-1}t}\mathbb R,
 \end{cases}
\]
以及来自 $\mathbb R^{d_h}$ 因子的常值拉格朗日道路
$t\mapsto\mathbb R$、边界条件
$\zeta(\partial[0,1]^2)\subset\mathbb R$ 的问题。对第一类问题，上述
$\partial[0,1]^2$ 上的回路为
\[
\begin{aligned}
 (s,0)&\longmapsto\mathbb R,&
 (0,t)&\longmapsto e^{-(l+1/2)\pi\sqrt{-1}t}\mathbb R,\\
 (s,1)&\longmapsto\sqrt{-1}\,\mathbb R,&
 (1,t)&\longmapsto e^{\pi\sqrt{-1}t/2}\mathbb R.
\end{aligned}
\]
沿 $\partial[0,1]^2$ 正向走过的这条回路，同伦于先走
\[
 t\longmapsto e^{(l+1/2)\pi\sqrt{-1}(1-t)}\mathbb R,
\]
再走
$t\longmapsto e^{\pi\sqrt{-1}t/2}\mathbb R$ 的拼接道路。该回路的
马斯洛夫指标为 $l+1$，求和得到
\[
 \mu(\widetilde\lambda_{0,\lambda_0,h})
 =
 \sum_{j=1}^{n-d_h}(l_j+1).
\]

另一方面，对应的弗雷德霍姆指标可用一复变量傅里叶分析显式求出；例如
\textbf{[Oh99]} 的 (A.12) 算得一维指标为 $l+1$。来自
$\mathbb C^{d_h}$ 的常值道路 $t\mapsto\mathbb R$ 虽有显然的常值解，
但它们不属于加权索伯列夫空间 $W_\delta^{1,p}$，因为其中元素必须由权重
$e^{\delta|\tau|}$ 强制指数衰减；相应余核也容易验证为零。因此
\[
 \operatorname{Index}
 \bar\partial_{\lambda_{w,\lambda_0},Z_-}
 =
 \sum_{j=1}^{n-d_h}(l_j+1).
\]
比较两个指标公式即得结论。
\end{proof}

下面讨论证明 \eqref{eq:3-7-60-4} 所需的定向平坦丛
$\Theta^-_{R_h}$；细节将在第 8.8 节给出。注意
\[
 \mathcal M_{k_1,k_0,(+\infty)}
 (L^{(1)},L^{(0)};[h,w],[h',w'];
 \vec P^{\,(1)},S,\vec P^{\,(0)})
\]
在仓西结构意义下未必可定向，这正是需要 $\Theta^-_{R_h}$ 的原因。
考虑纤维积
\[
\begin{aligned}
 &\mathcal M_{k_1,k_0}
 (L^{(1)},L^{(0)};[h,w],[h',w'])
 {}_{\operatorname{ev}}\mathop{\times}_{\prod_j f_j^{(1)}
 \times\prod_j f_j^{(0)}}
 \left(
 \prod_{j=1}^{k_1}P_j^{(1)}
 \times
 \prod_{j=1}^{k_0}P_j^{(0)}
 \right),
\end{aligned}
\]
记作
\[
 \mathcal M_{k_1,k_0}
 (L^{(1)},L^{(0)};[h,w],[h',w'];
 \vec P^{\,(1)},\vec P^{\,(0)}).
\]
弗洛尔理论中 $A_\infty$ 结构及边界算子的符号，由这个模空间的定向丛和
局部系统 $\{\Theta^-_{R_h}\}$ 决定。更精确地，置
$O_{R_{h'}}=\det TR_{h'}$。

\begin{lemma}{3.7.70}\label{lem:3-7-70}
空间
\[
 \mathcal M_{k_1,k_0}
 (L^{(1)},L^{(0)};[h,w],[h',w'];
 \vec P^{\,(1)},\vec P^{\,(0)})
\]
的仓西结构带有切丛。令 $U_p=V_p/\Gamma_p$ 与 $E_p$ 如仓西结构的定义
（定义 A1.1）中所述。映射
$\operatorname{ev}_{\pm\infty}$ 延拓为 $U_p$ 上的浸没；并且存在同构
\[
\begin{aligned}
 \det(E_p^*)\otimes\det(TV_p)
 \cong{}&
 \operatorname{ev}_{+\infty}^*\Theta^-_{R_{h'}}
 \otimes\operatorname{ev}_{+\infty}^*O_{R_{h'}}\\
 &\otimes\operatorname{ev}_{-\infty}^*\Theta^-_{R_h}.
\end{aligned}
\]
它与由切丛存在性得到的同构
\[
 \det(E_p^*)|_{V_{pq}}\otimes\det(TV_p)|_{V_{pq}}
 \cong
 \det(E_q^*)\otimes\det(TV_q),
 \qquad q\in\psi_p(s_p^{-1}(0)),
\]
相容。这里
\[
\begin{gathered}
 \operatorname{ev}_{-\infty}\colon
 \mathcal M_{k_1,k_0}
 (L^{(1)},L^{(0)};[h,w],[h',w'];
 \vec P^{\,(1)},\vec P^{\,(0)})\to R_h,\\
 \operatorname{ev}_{+\infty}\colon
 \mathcal M_{k_1,k_0}
 (L^{(1)},L^{(0)};[h,w],[h',w'];
 \vec P^{\,(1)},\vec P^{\,(0)})\to R_{h'}
\end{gathered}
\]
仍按前述方式定义。
\end{lemma}

注意 $\det(E_p^*)\otimes\det(TV_p)$ 正是上述仓西空间的定向丛。
又因 $M$ 已定向，作为 $\{\pm1\}$-平坦丛存在典范同构
\[
 \Theta^+_{R_{h'}}
 \cong O_{R_{h'}}\otimes\Theta^-_{R_{h'}}.
\]
详细论证将在第 8.8 节给出，见命题 8.8.6 与定义 8.8.11。那里研究
$\mathcal M(L^{(1)},L^{(0)};[h,w],[h',w'])$ 同带
$O_{R_h}\otimes\Theta^-_{R_h}$ 系数的奇异单形 $S$ 的纤维积，并在
定义 8.8.11 中给
\[
 \mathcal M_{k_1,k_0,(+\infty)}
 (L^{(1)},L^{(0)};[h,w],S,[h',w'])
\]
定向。借助交换因子时适当的符号规则，也可把这个模空间解释为前述模空间
同 $S$ 的纤维积，从而把结果表成引理 3.7.70 的形式。

\begin{proof}[命题 3.7.59 的证明（续）]
用引理 3.7.70 研究 \eqref{eq:3-7-60-4} 中的定向；细节见命题 8.8.7 的
证明。设
$(\phi\colon S\to R_h)\in\mathcal X_1(R_h)$，并在
$O_{R_h}\otimes\Theta^-_{R_h}\otimes O_S$ 上配备平坦截面
$s$，其中 $O_S=\det TS$。有
\[
\begin{aligned}
 \det(E_p^*)\otimes\det(T(V_p\times_{R_h}S))
 \cong{}&
 (-1)^{r_h}
 \operatorname{ev}_{+\infty}^*\Theta^+_{R_{h'}}
 \otimes\operatorname{ev}_{-\infty}^*\Theta^-_{R_h}\\
 &\otimes O_{R_h}\otimes O_S,
\end{aligned}
\]
其中 $r_h=\dim R_h$；因子 $(-1)^{r_h}$ 的来由也见命题 8.8.7 的证明。
因此
\[
\begin{aligned}
 &\operatorname{ev}_{+\infty}^*\Theta^+_{R_{h'}}
 \otimes\det(E_p^*)\otimes\det(T(V_p\times_{R_h}S))\\
 &\qquad\cong
 (-1)^{r_h(\mu([h,w])+1)}
 \operatorname{ev}_{-\infty}^*
 (O_{R_h}\otimes\Theta^-_{R_h})\otimes O_S.
\end{aligned}
\]
于是 $\phi^*(O_{R_h}\otimes\Theta^-_{R_h})\otimes O_S$ 的平坦截面 $s$
诱导
\[
 \operatorname{ev}_{+\infty}^*\Theta^+_{R_{h'}}
 \otimes\det(E_p^*)\otimes\det(T(V_p\times_{R_h}S))
\]
的平坦截面。这说明
\[
 \left(
 \mathfrak s_{[h,w],[h',w'],\vec P^{\,(1)},S,\vec P^{\,(0)}}^{-1}(0),
 \operatorname{ev}_{+\infty}
 \right)
\]
是 $R_{h'}$ 中以
$O_{R_{h'}}\otimes\Theta^-_{R_{h'}}
\cong\Theta^+_{R_{h'}}$ 为系数的链；这就证明了
\eqref{eq:3-7-60-4} 的定向部分。连同前述按第 7.2 节选择
$\mathcal X_1(R_h)$ 及相容多值截面的标准构造，命题 3.7.59 得证；留待后续
各章说明的事项除外。
\end{proof}

现在记
\[
\begin{aligned}
 &\mathcal M^{\mathfrak s}_{k_1,k_0,(+\infty)}
 (L^{(1)},L^{(0)};[h,w],[h',w'];
 \vec P^{\,(1)},S,\vec P^{\,(0)})\\
 &\qquad:=
 \left(
 \mathfrak s_{[h,w],[h',w'],\vec P^{\,(1)},S,\vec P^{\,(0)}}^{-1}(0),
 \operatorname{ev}_{+\infty}
 \right).
\end{aligned}
\]
其定向见定义 8.8.11。要把右端的虚基本链看成奇异链，须选取三角剖分以及
顶点次序，见定义 A1.28。

\begin{definition}{3.7.71}\label{def:3-7-71}
对 $S\in\mathcal X_1(R_h)$，以 $S\otimes[h,w]$ 表示
\[
 C(R_h;\Theta^-_{R_h})\otimes_{\mathbb Q}\Lambda_h
\]
中相应的带标签元素
\translatornote{原书称它是 $C(L^{(0)};\mathbb Q)$ 中的元素；但 $S$ 是
$R_h$ 上带 $\Theta^-_{R_h}$ 系数的链，而定义 3.7.50 又以 $[h,w]$ 生成
$\Lambda_h$，故该型别应是上式所示的带标签生成元直和项。}。令
\[
\begin{gathered}
 T^{\lambda_i^{(0)}}e^{\mu_i^{(0)}}P_i^{(0)}
 \in C(L^{(0)};\Lambda_{0,\mathrm{nov}}),\\
 T^{\lambda_i^{(1)}}e^{\mu_i^{(1)}}P_i^{(1)}
 \in C(L^{(1)};\Lambda_{0,\mathrm{nov}}),
\end{gathered}
\]
以及
\[
 T^\lambda e^\mu[S;h,w]
 \in C(R_h;\Theta^-_{R_h})\otimes\Lambda_{0,\mathrm{nov}}.
\]
定义
\[
\begin{aligned}
 &\mathfrak n_{k_1,k_0}
 \bigl(
 T^{\lambda_1^{(1)}}e^{\mu_1^{(1)}}P_1^{(1)}
 \otimes\cdots\otimes
 T^{\lambda_{k_1}^{(1)}}e^{\mu_{k_1}^{(1)}}P_{k_1}^{(1)}
 \otimes T^\lambda e^\mu(S\otimes[h,w])\\
 &\hspace{42mm}\otimes
 T^{\lambda_1^{(0)}}e^{\mu_1^{(0)}}P_1^{(0)}
 \otimes\cdots\otimes
 T^{\lambda_{k_0}^{(0)}}e^{\mu_{k_0}^{(0)}}P_{k_0}^{(0)}
 \bigr)\\
 &\quad=
 \sum_{[h',w']}
 T^{\lambda'}e^{\mu'}
 \mathcal M^{\mathfrak s}_{k_1,k_0,(+\infty)}
 (L^{(1)},L^{(0)};[h,w],[h',w'];
 \vec P^{\,(1)},S,\vec P^{\,(0)})
 \otimes[h',w'],
\end{aligned}
\]
其中
\[
 \lambda'
 =\sum_{i=1}^{k_0}\lambda_i^{(0)}
 +\lambda+\sum_{i=1}^{k_1}\lambda_i^{(1)},
 \qquad
 \mu'
 =\sum_{i=1}^{k_0}\mu_i^{(0)}
 +\mu+\sum_{i=1}^{k_1}\mu_i^{(1)}.
\]
当 $(k_1,k_0)\ne(0,0)$ 时采用上述定义；当
$(k_1,k_0)=(0,0)$ 时，还要加上
\[
 (-1)^{d_h+\mu([h,w])}\partial S\otimes[h,w].
\]
\end{definition}

定义 3.7.71 的求和包括
\[
 [h,w]=[h',w'],\qquad
 \lambda=\lambda',\qquad
 \mu=\mu',\qquad
 (k_1,k_0)\ne(0,0)
\]
的情形。对应项为
\[
\begin{aligned}
 \overline{\mathfrak m}_{k_1+1+k_0}
 \bigl(&P_1^{(1)}\cap R_h,\ldots,P_{k_1}^{(1)}\cap R_h,\,
 S,\\
 &P_1^{(0)}\cap R_h,\ldots,P_{k_0}^{(0)}\cap R_h\bigr),
\end{aligned}
\]
其中 $\overline{\mathfrak m}_{k_1+1+k_0}$ 是 $R_h$ 上杯积的一个
$A_\infty$ 形变。事实上，定义这类项的模空间由取值于 $R_h$ 的常值映射
组成。还可以选择 $L^{(0)},L^{(1)}$ 上的可数链集，使其都横截于 $R_h$；
否则上面的等式并不成立，不过模空间本身的横截性并不要求作这一选择。

又有
\[
 \deg S=d_h-\dim S,\qquad
 \deg P_j^{(i)}=n-\dim P_j^{(i)}.
\]
因此 \eqref{eq:3-7-60-1} 蕴含
\[
\begin{aligned}
 &d_{h'}-\dim
 \mathcal M^{\mathfrak s}_{k_1,k_0,(+\infty)}
 (L^{(1)},L^{(0)};[h,w],[h',w'];
 \vec P^{\,(1)},S,\vec P^{\,(0)})\\
 &\quad=
 \sum_j\deg'P_j^{(0)}
 +\sum_j\deg'P_j^{(1)}
 +\mu([h,w])-\mu([h',w'])+\deg S+1.
\end{aligned}
\]
故 $\mathfrak n_{k_1,k_0}$ 把移位次数提高 $1$。

\begin{theorem}{3.7.72}\label{thm:3-7-72}
上述 $\mathfrak n_{k_1,k_0}$ 在
$C(L^{(1)},L^{(0)};\Lambda_{0,\mathrm{nov}})$ 上定义了左
\[
 \bigl(C(L^{(1)};\Lambda_{0,\mathrm{nov}}),\mathfrak m^{(1)}\bigr)
\]
且右
\[
 \bigl(C(L^{(0)};\Lambda_{0,\mathrm{nov}}),\mathfrak m^{(0)}\bigr)
\]
的滤过 $A_\infty$ 双模结构。它是
$G(L^{(0)},L^{(1)})$-有隙的。此外，上述两个代数的同伦幺元对
$\{\mathbf e_1,\mathbf e_0\}$ 作为双模的同伦幺元作用。
\end{theorem}

\begin{proof}
这是第 3.7.4 小节中定理 3.7.21 之证明的直接推广，故略。
\end{proof}

现在假设两个 $A_\infty$ 代数
\[
 \bigl(C(L^{(0)};\Lambda_{0,\mathrm{nov}}),\mathfrak m^{(0)}\bigr),
 \qquad
 \bigl(C(L^{(1)};\Lambda_{0,\mathrm{nov}}),\mathfrak m^{(1)}\bigr)
\]
均无障碍。分别取其定界余链 $b_0,b_1$。由引理 3.7.14 与定理 3.7.72，
仿照第 3.7.4 小节得到上链复形
\begin{equation}
 \bigl(
 C(L^{(1)},L^{(0)};\Lambda_{0,\mathrm{nov}}),
 \delta_{b_1,b_0}
 \bigr).
\tag{3.7.73}\label{eq:3-7-73}
\end{equation}

\begin{definition}{3.7.74}\label{def:3-7-74}
把 \eqref{eq:3-7-73} 的上同调称为由 $b_0,b_1$ 形变的弗洛尔上同调，记作
\[
 HF((L^{(1)},b_1),(L^{(0)},b_0);\Lambda_{0,\mathrm{nov}}).
\]
\end{definition}
\liftIndex{I125}

当 $L^{(0)}$ 横截于 $L^{(1)}$ 时，复形 \eqref{eq:3-7-73} 与第 2.1.2
小节定义的复形一致。当 $L^{(0)}=L^{(1)}$ 时，有下述结论。

\begin{proposition}{3.7.75}\label{prop:3-7-75}
设 $L=L^{(0)}=L^{(1)}$，并且 $J_t$ 与 $t$ 无关。则定理 3.7.72 定义的
滤过 $A_\infty$ 双模，可以取成与例 3.7.6 所得双模相同。
\end{proposition}

\begin{proof}
此时 $L^{(1)}\cap L^{(0)}$ 只有一个分支，即 $L$。由 $L$ 与 $M$ 的
可定向性，局部系统 $\Theta^-_L$ 是平凡的。检视各项证明可知，可以取
$\mathcal X_1(R_h)=\mathcal X_1(L)$，其中 $\mathcal X_1(L)$ 如第 3.5 节，
而 $\mathcal X_1(R_h)$ 如命题 3.7.59。并且可取
\begin{equation}
\begin{aligned}
 &\mathcal M_{k_1,k_0,(+\infty)}
 (L^{(1)},L^{(0)};[h,w],[h,w'];
 \vec P^{\,(1)},S,\vec P^{\,(0)})\\
 &\qquad=
 \mathcal M_{k_1+k_0+1}
 (L;\beta;\vec P^{\,(1)},S,\vec P^{\,(0)}),
\end{aligned}
\tag{3.7.76}\label{eq:3-7-76}
\end{equation}
其中左端是命题 3.7.59 的构造，右端如命题 3.5.2，且
$w'=w\#\beta$。可以使命题 3.7.59 中的多值截面通过
\eqref{eq:3-7-76} 与命题 3.5.2 中的多值截面一致
\translatornote{原书先把 $\mathcal X_1(R_h)$ 误指向命题 3.7.55；又在
\eqref{eq:3-7-76} 后把左、右两端分别误指为命题 3.7.55、3.7.59，并把
待识别的两族多值截面都称为命题 3.5.2 中的多值截面。依本小节实际构造，
$\mathcal X_1(R_h)$、左端及其多值截面属于命题 3.7.59，右端及其多值
截面属于命题 3.5.2；译文据此恢复。}。于是由构造直接得到
\[
 \mathfrak n_{k_1,k_0}(x\otimes y\otimes z)
 =
 \mathfrak m_{k_1+k_0+1}(x,y,z).
\]
本书的定向约定与定义保证 \eqref{eq:3-7-76} 两端的定向确实一致，见
第 8.5、8.8 节。因此上述等式连同符号成立。
\end{proof}

由命题 3.7.75 得
\begin{equation}
 HF(L,b;\Lambda_{0,\mathrm{nov}})
 \cong
 HF((L,b),(L,b);\Lambda_{0,\mathrm{nov}}),
\tag{3.7.77}\label{eq:3-7-77}
\end{equation}
其中左端由滤过 $A_\infty$ 代数
$(C(L;\Lambda_{0,\mathrm{nov}}),\mathfrak m)$ 依定义 3.6.11 定义；在弱无
障碍情形则依定义 3.6.33 (1) 定义。右端依本小节的滤过
$A_\infty$ 双模定义。

再结合命题 3.7.75 与
\[
\begin{aligned}
 &HF((L^{(1)},b_1),(L^{(0)},b_0);\Lambda_{\mathrm{nov}})\\
 &\quad:=
 HF((L^{(1)},b_1),(L^{(0)},b_0);\Lambda_{0,\mathrm{nov}})
 \otimes_{\Lambda_{0,\mathrm{nov}}}\Lambda_{\mathrm{nov}}
\end{aligned}
\]
在哈密顿同痕下的不变性（见定理 5.3.14 与定理 4.1.5），便可利用第 6 章
建立的谱序列研究 $L$ 及其哈密顿形变的弗洛尔上同调。特别地，若
$L^{(1)}=\psi^1(L^{(0)})$ 是 $L^{(0)}$ 的哈密顿形变，将会看到
\begin{equation}
 HF(L,b;\Lambda_{\mathrm{nov}})
 \cong
 HF((\psi^1(L),\psi^1_*(b)),(L,b);\Lambda_{\mathrm{nov}}).
\tag{3.7.78}\label{eq:3-7-78}
\end{equation}
例子见下一小节。
